\section{Model Evaluation and Promotion}

\subsection{Model Evaluation}

% Describe how you evaluate the performance of your models.
% Include accuracy/error metrics, visual comparisons, or case studies.

\subsubsection{Advantages}

Summarize the strengths of your models (e.g., accuracy, robustness, interpretability, computational efficiency).

\subsubsection{Limitations}

Summarize the main limitations (e.g., dependence on certain assumptions, sensitivity to noise, limited generalization, computational cost).

\subsection{Future Work}

\subsubsection{Model Extension: Towards Game-Theoretic Robustness}

While our Adaptive Logistic Meritocracy System successfully addresses the fairness-satisfaction trade-off from a statistical perspective, future research should scrutinize its robustness through the lens of \textbf{Mechanism Design} and \textbf{Algorithmic Game Theory}.

\textbf{Incentive Compatibility (IC) and Strategic Voting:}
A critical extension is to formally verify whether the proposed system is Incentive Compatible (IC). In the current context, we assume fan votes are sincere expressions of preference. However, in a real-world adversarial setting, coordinated fan groups might engage in \textbf{Strategic Voting}—not only maximizing the score of their favorite contestant, but also conducting dishonest voting to increase the likelihood of their least favorite contestant being eliminated. Future work could model fans and judges as rational agents in a Bayesian Game, exploring whether the Adaptive System admits a Nash Equilibrium where truthful voting is the dominant strategy.

\textbf{Challenges and Data Limitations:}
The primary bottleneck for these economic extensions is data granularity. Our current model relies on aggregated vote shares ($\hat{\mathbf{S}}_t$). To rigorously test the excellent economic effects(e.g., IC condition), we require \textbf{individual-level ballot data} (e.g., identifying if user $i$ voted for contestant $A$ implies a zero probability of voting for rival $B$). The absence of such micro-data prevents us from observing collusion patterns or estimating the true utility functions of the voters.

\subsubsection{Model Application: Beyond the Show}

The core architecture of our model—a dual-track aggregation system with non-linear saturation—has broad applicability beyond reality television. It offers a generalizable framework for any domain requiring the integration of \textbf{"Expert Wisdom"} and \textbf{"Crowd Preference"}.

\textbf{Crowdsourcing and Peer Review:}
In platforms like Gitcoin Grants or academic peer review, where funding or acceptance decisions rely on community signals, "Sybil Attacks" and popularity contests are rampant risks. Our dynamic sensitivity parameter ($K$) could serve as an automated quality control filter: detecting anomalous variance in community votes and automatically increasing the weight of expert reviewers to stabilize the outcome.

\textbf{Real-Time Shadow Scoring System:}
For future seasons of \textit{Dancing with the Stars}, we envision deploying our model as a real-time \textbf{"Shadow Scorer"}. This system would run in parallel to the official results, providing producers with live metrics on "Fan Deviation" and "Fairness Risk." If the divergence between the official result and the Shadow Score exceeds a safety threshold, it could trigger an automatic "Judges' Choice" protocol, ensuring the competition remains mathematically grounded without human bias.